\documentclass[10pt,a4paper]{amsart}
\usepackage[margin=19mm]{geometry}
\usepackage{amsmath,amssymb,mathtools}
\usepackage[scr=rsfs,cal=euler]{mathalfa}
\usepackage{xfrac}
\usepackage[unicode,hidelinks,bookmarks=false]{hyperref}
\newcommand{\ie}{i.e.\ }
\newcommand{\cf}{cf.\ }
\newcommand{\resp}{respectively\ }
\newcommand{\Subsection}{\subsection}
\providecommand{\nobreakdash}{\nobreak\hbox{-}}
\setlength{\emergencystretch}{2em}
\multlinegap0pt
\usepackage{bm}
\usepackage{bbm}
\usepackage{dsfont}
\usepackage{xspace}
\usepackage{cancel}
\usepackage{mathtools}
\usepackage{amsthm}
\makeatletter
\@ifundefined{thm}{\newtheorem{thm}{Theorem}}{}
\@ifundefined{prop}{\newtheorem{prop}[thm]{Proposition}}{}
\makeatother
\newcommand{\ba}{\mathop{\mathrm{ba}}}
\newcommand{\ip}{\mathnormal{\mathsf{IP}}}
\newcommand{\sip}{\mathnormal{\mathsf{sIP}}}
\title[Finitely additive measures on $\mathbb Z$ and additive combi]{Finitely additive measures on $\mathbb Z$ and additive combinatorics}
\author{Zeinab Ashtab, David Fern\'andez-Bret\'on}
\date{Source: arXiv:2607.26522v1 (CC BY 4.0)}
\begin{document}
\maketitle

\noindent\emph{Source and licence.} Finitely additive measures on $\mathbb Z$ and additive combinatorics. Version arXiv:2607.26522v1 (CC BY 4.0), distributed under the Creative Commons Attribution 4.0 International licence, as recorded in the arXiv licence metadata for this exact version and shown on the abstract page https://arxiv.org/abs/2607.26522v1. Full rights notice: licenses/arxiv-2607.26522-CC-BY-4.0.txt.\par\smallskip

\noindent\textit{Editorial scope.} Counted unit: the Thm whose source label is \texttt{translationinvariant} in the article (source0.tex lines 704--706, source label \texttt{translationinvariant}), delivered together with the article's own complete original proof of it (source0.tex lines 707--809, one \texttt{proof} environment of 103 lines). No mathematics is written, repaired or re-derived: statement and proof are the article's own text line for line. Required context retained uncounted: characterizeipn (lines 371--373). Every definition, display and label of those lines is retained unchanged. Omitted from this package and not redistributed: 1--370, 374--703, 810--1049. Nothing inside a retained excerpt is omitted or altered: every retained body is delivered exactly as the article's lines are written, so no omission receipt of any kind accompanies this package. Citations: the retained excerpts carry no citation command at all, so no bibliography entry of the article is transcribed and no reference is redistributed. Reference adaptation: none was needed and none was made; every reference inside the delivered unit resolves inside it, so no unresolved reference or citation is produced. Printed slips kept literal: no printed word, symbol, hypothesis or number of the counted statement or of its proof was altered. Layout and class: the article's own class is \texttt{jloganal} with packages including ; the wrapper is the lane's pinned \texttt{amsart} editorial wrapper at 10pt on A4, which restates the theorem-like environments the retained text needs and adds the article's own macro definitions that the retained text needs (\texttt{\textbackslash ba}, \texttt{\textbackslash ip}, \texttt{\textbackslash sip}), with extra wrapper packages (bm, bbm, dsfont, xspace, cancel, mathtools). The delivered numbering is the unit's own. Assets: no retained span contains a figure environment, a graphics-inclusion command, a PDF-page inclusion, a table or a TikZ picture; no image file is copied into this package, the package declares no asset dependency and no image file is redistributed with it. Licence: version arXiv:2607.26522v1 of this article is distributed under the Creative Commons Attribution 4.0 International licence, as recorded in the arXiv licence metadata for this exact version and shown on the abstract page https://arxiv.org/abs/2607.26522v1. A full rights notice is delivered at licenses/arxiv-2607.26522-CC-BY-4.0.txt. Characters: every retained excerpt is pure ASCII, so the delivered engine, XeLaTeX with its default Latin Modern text font, reports no missing character.
\begin{prop}\label{characterizeipn}
Let $A\subseteq\mathbb Z$, and let $n\in\mathbb N$. Then, $A$ is an $\ip_{n+1}$-set if and only if there exists an $x\in A$ such that $A\cap(A-x)$ is an $\ip_{n}$-set.
\end{prop}

\begin{thm}\label{translationinvariant}
Let $\mu$ be a translation-invariant, positive finitely additive measure belonging to the unit ball of $\ba(\mathbb Z)$. Then, for every $A\subseteq\mathbb Z$, if $\mu(A)>1-\frac{1}{2^{n-1}}$ then $A$ is an $\sip_{n}$-set.
\end{thm}

\begin{proof}
We proceed by induction on $n$. The case $n=1$ is trivial since it obviously $\mu(A)>0$ implies $A\neq\varnothing$, that is, $A$ is $\sip_1$. Now suppose the result is true for $n$, and let $A\subseteq\mathbb Z$ be such that $\mu(A)>1-\frac{1}{2^n}$. Pick any $x\in A$ and consider the set $A-x$. Translation-invariance of $\mu$ implies that $\mu(A-x)=\mu(A)$, from where we can obtain
\begin{eqnarray*}
    1\geq \mu(A\cup(A-x)) & = & \mu(A)+\mu(A-x)-\mu(A\cap(A-x)) \\
    & = & 2\mu(A)-\mu(A\cap(A-x)) \\
    & > & 2-\frac{1}{2^{n-1}}-\mu(A\cap(A-x)).
\end{eqnarray*}
We may thus deduce that $\mu(A\cap(A-x))>1-\frac{1}{2^{n-1}}$; by induction hypothesis, this implies that $A\cap(A-x)$ is an $\sip_n$-set, and so $A$ must be an $\sip_{n+1}$-set by Lemma~\ref{characterizeipn}.
%To show $A$ is $IP_{n}$, we must produce numbers
%\[
%x_{1}, x_{2}, \dots, x_{n}
%\]
%such that all sums of at most $n$ distinct $x_{i}$'s lie in $A$, for every nonempty $F\subseteq\{1.\dots,n\}$ is following,
%\[
%\sum _{i\in F}\, x_{i}\in A,
%\]
% we choose them inductively, and at step $K+1$ we are forced to look at an intersection of several shifts of $A$,
% \[
% \bigcap_{F\subseteq\{1,\dots,k\}} (A-\sum_{i\in F}\, x_{i}).
% \]
% To be able to pick $x_{k+1}$ from that intersection we need its $\mu$-measure to be greater than zero.\\
% Now we should prove that the intersection of finitely many such shifts still have positive measure.\\
% $\mu(A)>0$, so we can choose $x_{1}\in A$. Suppose we have already chosen $x_{1},\dots,x_{k}$ $(k<n)$ such that new sums (those involving $x_{k+1})$ lie in $A$. That is of the form 
% \[
% x_{k+1}+s, \,\,\, s\in FS^{\leq k}(x_{1},\dots,x_{k}).
% \]
% For such a sum to be in $A$, we need $x_{k+1} \in A-s$. We require $x_{k+1} \in \cap_{s\in FS^{\leq k}}(A-s)$.\\
% We denote $E_{k}:=\bigcap_{s\in FS^{\leq k}}(A-s)$, and we claim that $\mu(E_{k})>0$.\\
% Each set $A-s$ is a shift of $A$, so 
% \[
% \mu(A-s)=\mu(A)>1-\frac{1}{2^{n}}.
% \]
% The set $FS^{(k)}$ has at most $2^{k}-1$ nonempty sums, so $2^{k}\leq 2^{n-1}<2^{n}$, $(k<n)$.
% \begin{align}\label{sum1}
%     \mu\Big(\cup_{s}\big(\mathbb N\setminus(A-s)\big)\Big)\leq\sum_{s}\mu\Big(\mathbb N\setminus(A-s)\Big)=|s|\Big(1-\mu(A)\Big)
% \end{align}
% Dr. Breton, The first inequality is because of for any finite family of sets, and the second equality is since of complements:
% \begin{align*}
%   &\mu(\mathbb N\setminus C)=1-\mu(C) \\&
%   \mu(\mathbb N\setminus(A-s))=1-\mu(A-s)=1-\mu(A).
% \end{align*}

%We know that $1-\mu(A)<\frac{1}{2^{n}}$, so from (\ref{sum1})
%\[
% \mu\Big(\cup_{s}\big(\mathbb N\setminus(A-s)\big)\Big)< 2^{k}. \frac{1}{2^{n}}\leq 2^{n-1}. \frac{1}{2^{n}}=\frac{1}{2}<1.
%\]
%Therefore the complement of the union satisfies 
%\[
%\mu(E_{k})=1-\mu\Big(\cup_{s}\big(\mathbb N\setminus(A-s)\big)\Big)>0,
%\]
%so $E_{k}$ is nonempty andd finally we can choose $x_{k+1}\in E_{k}$, and 
%\[
%FS^{\leq k+1}(x_{1},\dots,x_{k+1})\subseteq A.
%\]
%At the end we can get numbers $x_{1},\dots,x_{n}$ such that 
%\[
%FS^{\leq n}(x_{1},\dots,x_{n})\subseteq A,
%\]
%this means $A$ is an $IP_{n}$-set.
%\qed
%\proof (2)
%Let $x\in A$ such that $ \mu(A-x)=\mu(A)$ and suppose $\mu (A)>1-\frac{1}{2^{n+1}}$ we have,
%\[
%\mu(A\cup(A-x))=\mu(A)+\mu(A-x)-\mu(A\cap(A-x)).
%\]
%Thus
%\[
%\mu(A\cap(A-x))\geq 1-(1-\frac{1}{2^{n+1}})-(1-\frac{1}{2^{n+1}})=1-\frac{1}{2^{n}}.
%\]
%By induction hypothesis $ A\cap (A-x)$ is an $IP_{n}$, and by the previous proposition %(\ref{removingxofA,isIP}) $A$ is $IP_{n+1}$.
%(3) For $n=1$, define $\alpha_{1}=0$, if $\mu(A)>0$, then $A\neq \varnothing$ hence $A$ is $IP_{1}$-set. So the statement holds for $n=1$.\\
%Define $\alpha_{n+1}=\frac{1+\alpha_{n}}{2}$, this gives an increasing sequence 
%\[
%0=\alpha_{1}<\alpha_{2}=\frac{1}{2}<\alpha_{3}=3
%\frac{3}{4}<\dots<1,
%\]
%with
%\[
%\alpha_{n}=1-\frac{1}{2^{n-1}}.
%\]
%Idempotence gives translation invariance (in the averaged sense). $\mu$ is idempotent so
%\[
%\mu(A)=\mu(n\to \mu(A-n)).
%\]
%Let $B=\{n\in S \,|\, \mu(A\cap(A-n))>\alpha_{n}\}$. We claim $\mu(B)>0$. If $\mu(B)=0$, %then for $\mu$-almost all $n$,
%\[
%\mu(A\cap(A-n))\leq\alpha_{n},
%\]
%but
%\[
%\mu(A\cap(A-n))=\mu(A)+\mu(A-n)-\mu(A\cup(A-n))\leq 2\mu(A)-\mu(A\cup(A-n)),
%\]
%by averaging and using idempotence,
%\[
%\mu(A)\leq\alpha_{n}+(1-\mu(A))
%\]
%which implies that
%\[
%\mu(A)\leq \frac{\alpha_{n}+1}{2}=\alpha_{n+1}
%\]
%this contradicting $\mu(A)>\alpha_{n+1}$, hence $\mu(B)>0$, so by induction hypothesis %$B$ is $IP_{n}$. Therefore there exist $x_{1},\dots,x_{n}\in B$ such that %$FS(x_{1},\dots,x_{n})\subset B$.\\
%Add one more element and pick $x_{n+1}\in A\cap\bigcap_{s\in FS(x_{1},\dots,x_{n})} (A-%s)$. This intersection is nonempty because each term has measure greater than %$\alpha_{n}$, and finitely many such intersections preserve positivity. Then %$FS(x_{1},\dots,x_{n+1)}\subset A$, hence $A$ is $IP_{n+1}.$
\end{proof}

\end{document}
